\chapter{Reproducibility Configuration and Interfaces}
\label{app:reproducibility}

This appendix identifies the repository contracts required to reproduce the
accepted analyses. It is an index to the preserved source and evidence rather
than a replacement for the raw bundles. Repository-relative paths are given so
that a result can be traced without copying large logs or model artifacts into
the manuscript. Every rerun should use a new output directory and must not
overwrite the source evidence.

\section{Evidence and configuration hierarchy}

The reproduction hierarchy has four layers. Raw device, training and simulator
outputs are the primary observations. Machine-readable manifests bind those
outputs to file identities and hashes. Analysis scripts derive tables and
figures from the manifests. The thesis reports the accepted values and their
interpretation boundaries. When these layers disagree, the raw artifact and
its contemporaneous manifest take precedence over prose summaries.

\begin{table}[H]
  \centering
  \caption{Principal frozen contracts used by the integrated thesis.}
  \label{tab:appendix-contracts}
  \scriptsize
  \begin{tabularx}{\textwidth}{@{}p{0.30\textwidth}p{0.23\textwidth}X@{}}
    \toprule
    Repository-relative path & Role & Status represented in this thesis \\
    \midrule
    \path{manifests/fire_only_curated_v7.freeze.json} & Dataset identity,
    counts and audit hashes & Frozen fire-only source-family split used for the
    final six-model study. \\
    \path{config/six_model_validation_v1.json} & Common validation contract &
    Frozen 700-image validation evaluation and threshold-selection rule. \\
    \path{config/model_deployment_registry_v1.json} & Native model registry &
    Six screened native artifacts; the strengthened platform matrix was not
    executed. \\
    \path{config/exact_device_screening_v1.json} & Device-screen protocol &
    Registered 20-warm-up/200-measurement protocol; not completed as a six by
    two matrix. \\
    \path{ros2_ws/src/drone_enhanced_nodes/config/runtime.yaml} & Node defaults
    and safety controls & Versioned defaults only; executed trials retain their
    own launch parameters and logs. \\
    \path{config/formal_simulation_campaign_v1.json} & Formal simulation
    campaign & Preregistered campaign; development pilots are reported
    separately because the complete campaign was not executed. \\
    \path{config/external_evaluation_sources_v1.json} & External-source gates &
    Readiness and eligibility contract; no external inference result was
    admitted. \\
    \path{evidence/manifests/evidence_manifest_20260804.csv} & Evidence
    inventory & File size, UTC modification time and SHA-256 inventory for the
    preserved evidence snapshot. \\
    \bottomrule
  \end{tabularx}
\end{table}

The frozen dataset contains 4,991 records: 3,558 training, 700 validation and
733 test images. The record manifest SHA-256 is
\path{478803efb7f8f4f6c3fe74d58658958bccfa2d0f58ec50a1a5c16535e393a7dd}.
The validation COCO annotation SHA-256 recorded in the common contract is
\path{1a93f800d63bf355bb586039da245cd612cbdd7052102cbfd6da1409312ab7c6}.
These identities distinguish the final curated experiment from the earlier
historical Sentinel and Scout datasets.

\section{Reproduction stages and entry points}

A complete rerun proceeds in the following order. Each stage consumes the
preserved output of the preceding stage and writes a new derived bundle.

\begin{enumerate}
  \item \textbf{Validate the contracts.} Use
  \path{tools/validate_strengthening_config.py},
  \path{tools/validate_screening_contract.py} and
  \path{tools/validate_deployment_and_external_contracts.py} before training or
  evaluation.
  \item \textbf{Materialize the frozen data.} Dataset construction and export
  are implemented by \path{tools/build_curated_research_manifest.py},
  \path{tools/materialize_training_dataset.py},
  \path{tools/export_training_coco.py} and
  \path{tools/export_training_tfrecord.py}. The generated counts and hashes
  must match the frozen manifest.
  \item \textbf{Run model evaluation.} The framework-specific prediction entry
  points are \path{tools/predict_yolo_coco.py},
  \path{tools/predict_tfod_coco.py} and
  \path{tools/predict_tflite_coco.py}. Common metric calculation uses
  \path{tools/evaluate_coco_predictions.py}; grouped F1 intervals use
  \path{tools/grouped_bootstrap_detection.py}. Thresholds are chosen on the
  validation split and must not be retuned on the opened test split.
  \item \textbf{Validate conversion fidelity.} Candidate generation and checks
  use \path{tools/export_yolo_onnx_candidates.py},
  \path{tools/export_ssd_tflite_candidate.py},
  \path{tools/validate_yolo_onnx_fidelity.py} and
  \path{tools/validate_prediction_fidelity.py}. A converted artifact that
  fails the frozen fidelity rule is excluded from equivalent-model device
  ranking.
  \item \textbf{Analyze retained device evidence.} Backend, communication and
  integrated-workflow bundles are processed by
  \path{tools/analyze_backend_benchmarks.py},
  \path{tools/analyze_ros_roundtrip.py} and
  \path{tools/analyze_integrated_trial.py}. These tools operate on completed
  evidence; they do not turn an incomplete run into an accepted result.
  \item \textbf{Check thesis transcription.} The numerical cross-check entry
  point is \path{tools/verify_thesis_numeric_contract.py}. The manuscript is
  then compiled with pdfLaTeX and Biber, and its pages are inspected visually.
\end{enumerate}

The scripts expose their current arguments in the repository and should be run
from the repository root in the environment recorded for the corresponding
experiment. Environment locks are stored under \path{config/} and
\path{simulation/}. This appendix intentionally does not prescribe a single
global Python environment because Ultralytics, TensorFlow Object Detection API,
EfficientDet and ROS~2 were executed in distinct compatible environments.

\section{Deployment artifact identities}

Table~\ref{tab:appendix-artifacts} records the six native candidates used by
the deployment-preparation register. A native training artifact and a converted
runtime artifact are different identities; neither may inherit the other's
performance result without a recorded conversion and fidelity check.

\begin{table}[H]
  \centering
  \caption{Native candidate identities from the deployment registry.}
  \label{tab:appendix-artifacts}
  \scriptsize
  \begin{tabularx}{\textwidth}{@{}p{0.23\textwidth}p{0.12\textwidth}X@{}}
    \toprule
    Model ID & Input & Native artifact SHA-256 \\
    \midrule
    \path{yolov8n_416} & 416 &
    \path{854482d0fbd320e51575f0141a479e517d721ff6cd512d12c60d0acf30a35ff3} \\
    \path{yolov8n_640} & 640 &
    \path{1c49f90ee7e635b9011ae1ec34a9cfccfea75f9b08941d7475d031587a0c0a33} \\
    \path{yolov8s_416} & 416 &
    \path{7bf6137add173f51d015c18841b5903300db909d31fb82ea3f51f3498ee366c5} \\
    \path{yolov8s_640} & 640 &
    \path{53019aa916dcbfa11f3f38a3cc16579270436e44fa7bd80ccd1fc37afb60f0f0} \\
    \path{ssd_mobilenet_v2_fpnlite_320} & 320 &
    \path{ee20414e69c0473e17989f19724709708c9a49303e2e26e8eabe4d771cf9e229}
    (checkpoint index; the registry also records the data shard). \\
    \path{efficientdet_lite0_320} & 320 &
    \path{948e71fa53fa9a740a2bca45f169280a9c6e6596a49bfcab65d0ad9cfda2f12c} \\
    \bottomrule
  \end{tabularx}
\end{table}

The final deployed Sentinel engine is the 640-pixel TensorRT artifact with
input shape $[1,3,640,640]$ and output shape $[1,6,8400]$. The historical
416-pixel engine has output shape $[1,6,3549]$. The two identities remain
separate throughout Chapters~\ref{ch:implementation}--\ref{ch:conclusion}; the
older latency measurements are never assigned to the final engine.

\section{Correlated ROS 2 interfaces}

The enhanced message definitions are stored under
\path{ros2_ws/src/drone_enhanced_interfaces/msg/}. Table
\ref{tab:appendix-interfaces} groups the fields by analytical purpose. Field
names and types remain defined by the message files, which are the authoritative
interface source.

\begin{table}[H]
  \centering
  \caption{Analytical purpose of the enhanced ROS~2 message fields.}
  \label{tab:appendix-interfaces}
  \scriptsize
  \begin{tabularx}{\textwidth}{@{}p{0.21\textwidth}p{0.31\textwidth}X@{}}
    \toprule
    Message & Principal fields & Purpose \\
    \midrule
    \texttt{HazardAlert} & \texttt{experiment\_id}, \texttt{trial\_id},
    \texttt{event\_id}, hazard type, confidence, temporal count, frame/alert
    monotonic times, model hash and normalized target box & Identifies the
    Sentinel request and preserves the detection context. \\
    \texttt{HazardVerdict} & matching IDs, verdict, hazard type, Scout
    confidence, receipt/frame/verdict monotonic times, model hash and reason &
    Returns a correlated Scout outcome without claiming independent physical
    ground truth. \\
    \texttt{HazardDecision} & matching IDs, final decision, both confidences,
    detection/alert/verdict/decision monotonic times and reason & Records the
    terminal system decision and supports latency decomposition. \\
    \bottomrule
  \end{tabularx}
\end{table}

Both physical nodes retained \texttt{enable\_flight\_commands=false} in the
evaluated stationary workflow. Image-space target coordinates were logged for
traceability; the Scout used its latest full frame and did not navigate to or
crop around the transferred box.

\section{Platform integration and environment inventory}

Table~\ref{tab:appendix-physical-integration} records the physical integration
facts admitted to the thesis and prevents an undocumented installation detail
from becoming a quantitative claim. The topology is reproducible at the
component level; the missing electrical and geometric measurements remain
explicit.

\begin{table}[H]
  \centering
  \caption{Physical integration record and measurement boundary.}
  \label{tab:appendix-physical-integration}
  \scriptsize
  \begin{tabularx}{\textwidth}{@{}p{0.22\textwidth}p{0.39\textwidth}X@{}}
    \toprule
    Item & Retained fact & Unmeasured or unverified detail \\
    \midrule
    Sentinel airframe & Customized Holybro X500 V2 with Pixhawk~6C & Installed
    mass, dimensions, rotor spacing, centre of gravity and payload margin \\
    Companion power & Jetson Nano supplied from airframe power distribution
    through an intermediate DC--DC regulator & Regulator identity, input/output
    time series, efficiency and energy per inference \\
    Flight-controller link & Three-conductor physical link; retained service
    uses \path{/dev/ttyTHS1} at 57,600~baud & Connector-level pin map and
    electrical signal verification \\
    Sentinel camera & D435i attached to the Sentinel and connected to the
    Jetson; colour endpoint exercised & Camera-to-body translation/rotation,
    depth use and vibration response \\
    Scout platform & ModalAI development platform with integrated VOXL~2,
    camera and PX4 services & Independently measured mass, geometry and
    flight-performance parameters \\
    Physical safety state & Stationary aircraft with propellers removed and
    flight commands disabled & No physical integrated flight result \\
    \bottomrule
  \end{tabularx}
\end{table}

Table~\ref{tab:appendix-environment-inventory} freezes the software-generation
facts used to interpret the deployment. An unavailable version is stated as
such rather than reconstructed from a later vendor image.

\begin{table}[H]
  \centering
  \caption{Onboard operating-environment inventory.}
  \label{tab:appendix-environment-inventory}
  \scriptsize
  \begin{tabularx}{\textwidth}{@{}p{0.22\textwidth}p{0.34\textwidth}X@{}}
    \toprule
    Layer & Sentinel & Scout \\
    \midrule
    Compute & Jetson Nano 4~GB & VOXL~2 / QRB5165 \\
    Host OS & Ubuntu 18.04.6 on Jetson Linux R32.7.6 / JetPack 4.6.6,
    kernel \texttt{4.9.337-tegra} & Exact base distribution/version not
    retained in the accepted inventory \\
    ROS~2 & Foxy inside
    \path{dustynv/ros:foxy-ros-base-l4t-r32.7.1} & Native Foxy at
    \path{/opt/ros/foxy} \\
    Build/runtime boundary & Docker 20.10.21, NVIDIA runtime, host networking,
    bind-mounted workspace & Native \texttt{colcon} workspace and vendor
    services \\
    Perception dependencies & D435i, GStreamer, CUDA, TensorRT and Python node &
    MPA camera server/bridge, YUV/RGB decoder, TensorFlow Lite and Python node \\
    DDS configuration & \multicolumn{2}{p{0.73\textwidth}@{}}{Cyclone DDS; SwarmNet
    \texttt{wlan0}; multicast; isolated domain per experiment; matching enhanced
    interfaces} \\
    \bottomrule
  \end{tabularx}
\end{table}

Table~\ref{tab:appendix-node-deployment} maps the principal processes to their
host and communication role. Helper probes used in individual experiments are
excluded because they are measurement tools rather than required members of
the deployed hazard-decision path.

\begin{table}[H]
  \centering
  \caption{Principal onboard node and service deployment.}
  \label{tab:appendix-node-deployment}
  \scriptsize
  \begin{tabularx}{\textwidth}{@{}p{0.20\textwidth}p{0.17\textwidth}p{0.24\textwidth}X@{}}
    \toprule
    Process or service & Host & Publishes & Subscribes / dependency \\
    \midrule
    Sentinel perception/coordinator & Jetson Foxy container &
    \path{/swarm/hazard_alert}, \path{/swarm/hazard_decision} & D435i frame,
    \path{/swarm/hazard_verdict}, TensorRT engine \\
    Scout verifier & VOXL~2 native Foxy & \path{/swarm/hazard_verdict} &
    \path{/swarm/hazard_alert}, latest MPA image, TFLite model \\
    MPA-to-ROS~2 bridge & VOXL~2 vendor environment & Scout image topic & ModalAI
    camera pipe and camera-server service \\
    MAVProxy telemetry service & Jetson host & Laboratory telemetry forwarding &
    Pixhawk serial endpoint; not a perception-decision dependency \\
    Asynchronous evidence writers & Respective edge node & Device-local JSONL
    and manifest files & Local node-event queues \\
    SwarmNet access point & Sentinel host & Local wireless connectivity &
    Physical radio and access-point service \\
    Laboratory laptop & External observer & Test commands and retained copies &
    Not required for onboard inference or final decision \\
    \bottomrule
  \end{tabularx}
\end{table}

\section{Reproduction limits}

Reproducibility does not remove the validity limits reported in
Chapter~\ref{ch:discussion}. The device operating systems, drivers and hardware
state can affect latency; training used one controlled seed per configuration;
and the physical trials contain dependent requests within stationary
exposures. Repeating a script can verify the recorded computation, but it does
not make the sample representative of industrial conditions.
